% !PRINTED 312
\weberpage{312}
%% [page furniture, not body text -- commented out 2026-08-10] Page furniture (running
%% head, folio, signature) is excluded from this edition by standing policy.
%% \noindent Zehnter Abschnitt. \hfill \S.\,81.

Wenn wir in der Summe $\Sigma \alpha\xi$ für das System $(\xi)$ alle sieben möglichen Annahmen machen, so erhält man dreimal den Werth 0 und folglich viermal den Werth 1.

Die Congruenz
\weq{7}{\Sigma \alpha\xi \equiv 1 \ (\mathrm{mod}\ 2)}
hat also vier und nur vier verschiedene Lösungen für $\xi$. Die Gesammtheit dieser vier Lösungen wollen wir ein \gesperrt{Quadrupel} nennen. Jedes Tripel bestimmt eindeutig ein Quadrupel und umgekehrt, und es giebt also auch sieben verschiedene Quadrupel, die nach (7) durch das Zahlensystem $(\alpha)$ vollständig bestimmt sind. Wir beweisen nun den Satz:

\gesperrt{6. Durchläuft $(\xi)$ das zu $(\alpha)$ gehörige Quadrupel und ist $(\beta)$ ein von $(\alpha)$ verschiedenes System, so sind unter den vier Summen $\Sigma \beta\xi$ zwei gerade und zwei ungerade.}

Denn wenn $(\xi)$ die Gesammtheit aller sieben Systeme durchläuft, so finden sich unter den Summen $\Sigma \beta\xi$ drei gerade und vier ungerade; von diesen liefert das zu $(\alpha)$ gehörige Tripel zwei ungerade und eine gerade; es bleiben also für das Quadrupel zwei gerade und zwei ungerade.

Wir ändern nun die Bezeichnung bei den Quadrupeln und nehmen für das System $(\alpha)$ ein anderes Zeichen $(h_1, h_2, h_3) = (h)$. Das zu $(h)$ gehörige Quadrupel bezeichnen wir mit $(\alpha), (\beta), (\gamma), (\delta)$, so dass die vier Congruenzen:
\weq{8}{\Sigma h\alpha \equiv 1, \quad \Sigma h\beta \equiv 1, \quad \Sigma h\gamma \equiv 1, \quad \Sigma h\delta \equiv 1 \quad (\mathrm{mod}\ 2)}
bestehen.

Im Nenner des Productes der vier Functionen
\[
v_\alpha, \quad v_\beta, \quad v_\gamma, \quad v_\delta
\]
kommen nach (6) im Ganzen 12 Factoren $\Psi_\xi$ vor, und zwar kommt jeder Factor $\Psi_\xi$ so oft darunter vor, als die Anzahl der geraden Zahlen unter den Summen
\[
\Sigma \xi\alpha, \quad \Sigma \xi\beta, \quad \Sigma \xi\gamma, \quad \Sigma \xi\delta
\]
beträgt, d.\,h.\ $\Psi_h$ kommt gar nicht vor, während jedes andere $\Psi_\xi$ zweimal vorkommt (nach 6.).

Wenn wir also noch die Relation (1) benutzen, so ergiebt sich:
\weq{9}{\Psi_h^2 = A^2 v_\alpha v_\beta v_\gamma v_\delta.}
